\documentclass[a4paper,10.5pt]{article}
\usepackage[T1]{fontenc}\usepackage[utf8]{inputenc}
\usepackage[top=2.2cm,bottom=2.2cm,left=2.35cm,right=2.35cm]{geometry}
\usepackage{hyperref}\usepackage{xcolor}\definecolor{accent}{HTML}{1F4E78}\hypersetup{colorlinks=true,urlcolor=accent}
\pagestyle{empty}\setlength{\parindent}{0pt}\setlength{\parskip}{7pt}
\begin{document}
\textbf{Prabhavit Gupta}\\
Venice, Italy \textbar{} +39 351 610 7807 \textbar{} prabhavitg@gmail.com\\[10pt]
Dear hello again Hiring Team,

I am writing to apply for the AI Solutions Engineer (d/f/m) role. I have three years of continuous professional experience across eight European companies, working at the point where a business use case turns into a working AI pipeline, and I am comfortable moving between stakeholder conversations and hands on prototyping in the same week.

At UltraStrategy, I built the CV parsing pipeline behind the platform: a user uploads a CV, the pipeline extracts and structures the data with Affinda and the Claude API, scores how relevant the candidate is for specific types of roles, and the result is delivered to the frontend before a standardised profile is generated. That is the same shape of work this role asks for, translating a concrete business need into a working data pipeline, evaluating output quality, and shipping it into a live product, and I also built the platform serving the different user profiles behind it.

At FlowSpace, a live workspace marketplace, I built support AI agents that answer customer queries from live platform data and escalate to a human on low confidence, and I own an AI tool that turns any brochure or web page into a structured listing with a human in the loop, over 900 produced so far. I also built the CRM and GTM automations behind lead capture, enrichment and outbound. Across both, the pattern is the same one you describe: identify where AI adds measurable business value, prototype fast, and put it into production with the right guardrails.

I write Python and TypeScript, work directly with LLM APIs including the Claude API and OpenAI, and have hands on experience with RAG grounding and n8n based data pipelines. I hold Italian EU long term resident status. I am not currently authorised to work in Austria and would need the applicable employer-side permit or sponsorship process, and I am ready to relocate to Leonding or Vienna and complete that process.

I would welcome the opportunity to discuss how I could contribute to hello again's Intelligence Layer.

Regards,\\
Prabhavit Gupta
\end{document}
